\documentclass[10pt, aspectratio=169]{beamer}
% Preámbulo modular para presentaciones Beamer (Modelación Estadística)
% Diseño y paleta institucional del Tecnológico de Monterrey

\documentclass[aspectratio=169,xcolor={dvipsnames,table}]{beamer}

\usepackage[utf8]{inputenc}
\usepackage[spanish,mexico]{babel}
\usepackage{amsmath,amssymb,amsfonts,mathtools}
\usepackage{graphicx}
\usepackage{listings}
\usepackage{booktabs}
\usepackage{tikz}
\usetikzlibrary{matrix,arrows,decorations.pathmorphing,shapes.geometric,positioning}

% Paleta Institucional Tecnológico de Monterrey
\definecolor{TecAzul}{HTML}{0039A6}       % Azul TEC: color primario institucional
\definecolor{TecAzulOscuro}{HTML}{1A2E51} % Azul marino oscuro
\definecolor{TecRojo}{HTML}{EC2661}       % Rojo de acento (paleta miTec)
\definecolor{TecGris}{HTML}{646464}       % Gris neutro para comentarios y subtítulos
\definecolor{TecFondoBloque}{HTML}{F4F6F9}% Fondo suave para bloques

% Configuración de Tema Beamer
\usetheme{Boadilla}
\usecolortheme[named=TecAzulOscuro]{structure}
\setbeamercolor{alerted text}{fg=TecRojo}
\setbeamercolor{palette primary}{bg=TecAzulOscuro,fg=white}
\setbeamercolor{palette secondary}{bg=TecAzul,fg=white}
\setbeamercolor{palette tertiary}{bg=TecAzulOscuro!80!black,fg=white}
\setbeamercolor{title}{fg=TecAzulOscuro,bg=TecFondoBloque}
\setbeamercolor{frametitle}{fg=TecAzulOscuro,bg=TecFondoBloque}

% Bloques personalizados elegantes
\setbeamercolor{block title}{bg=TecAzulOscuro,fg=white}
\setbeamercolor{block body}{bg=TecFondoBloque,fg=black}
\setbeamercolor{block title alerted}{bg=TecRojo,fg=white}
\setbeamercolor{block body alerted}{bg=TecRojo!8,fg=black}
\setbeamercolor{block title example}{bg=TecAzul,fg=white}
\setbeamercolor{block body example}{bg=TecAzul!8,fg=black}

% Redondear bordes y añadir sombra sutil
\setbeamertemplate{blocks}[rounded][shadow=true]
\setbeamertemplate{navigation symbols}{}

% Pie de página limpio con número de diapositiva
\setbeamertemplate{footline}{
  \leavevmode%
  \hbox{%
  \begin{beamercolorbox}[wd=.7\paperwidth,ht=2.25ex,dp=1ex,left]{author in head/foot}%
    \hspace*{2ex}\insertshortauthor~~(\insertshortinstitute) \hfill \insertshorttitle
  \end{beamercolorbox}%
  \begin{beamercolorbox}[wd=.3\paperwidth,ht=2.25ex,dp=1ex,right]{date in head/foot}%
    \insertshortdate{}\hspace*{2em}
    \insertframenumber{} / \inserttotalframenumber\hspace*{2ex}
  \end{beamercolorbox}}%
  \vskip0pt%
}

% Configuración de Listings de Python para visualización pedagógica de código
\lstset{
    language=Python,
    basicstyle=\ttfamily\footnotesize,
    keywordstyle=\color{TecAzulOscuro}\bfseries,
    stringstyle=\color{TecRojo},
    commentstyle=\color{TecGris}\itshape,
    showstringspaces=false,
    breaklines=true,
    frame=lines,
    rulecolor=\color{TecAzulOscuro!30},
    backgroundcolor=\color{TecFondoBloque},
    aboveskip=0.5em,
    belowskip=0.5em,
    literate={á}{{\'a}}1 {é}{{\'e}}1 {í}{{\'i}}1 {ó}{{\'o}}1 {ú}{{\'u}}1
             {Á}{{\'A}}1 {É}{{\'E}}1 {Í}{{\'I}}1 {Ó}{{\'O}}1 {Ú}{{\'U}}1
             {ñ}{{\~n}}1 {Ñ}{{\~N}}1 {¿}{{?`}}1 {¡}{{!`}}1
}

% Metadatos institucionales por defecto
\title[Modelación Estadística]{Modelación Estadística para Ciencia de Datos e Ingeniería}
\author[J. Castillo Colmenares]{Juliho David Castillo Colmenares}
\institute[Tec de Monterrey]{Tecnológico de Monterrey \\ Escuela de Ingeniería y Ciencias}
\date{\today}

% Comandos y macros estadísticas compartidas para las presentaciones Beamer (Metropolis)

% Conjuntos numéricos básicos
\newcommand{\R}{\mathbb{R}}
\newcommand{\N}{\mathbb{N}}
\newcommand{\Q}{\mathbb{Q}}
\newcommand{\Z}{\mathbb{Z}}
\newcommand{\set}[1]{\left\{ #1 \right\}}
\newcommand{\abs}[1]{\left|#1\right|}
\newcommand{\norm}[1]{\left\| #1 \right\|}

% Comandos estadísticos y probabilísticos del curso
\newcommand{\Var}{\operatorname{Var}}
\newcommand{\cov}{\operatorname{Cov}}
\newcommand{\corr}{\rho}
\newcommand{\comb}[2]{\begin{pmatrix} #1 \\ #2 \end{pmatrix}}
\newcommand{\s}{\sigma}
\newcommand{\particion}{\mathfrak{P}}
\newcommand{\card}[1]{\#\left( #1 \right)}
\newcommand{\seq}[3]{\set{#1}_{#2}^{#3}}
\newcommand{\E}{\mathbb{E}}
\newcommand{\Prob}{\mathbb{P}}

% Entornos pedagógicos y cajas en español académico natural
\newenvironment{discusion}{%
  \begin{alertblock}{Pregunta de Discusión}%
}{%
  \end{alertblock}%
}

\newenvironment{reflexion}{%
  \begin{alertblock}{Reflexión para el Aula}%
}{%
  \end{alertblock}%
}

\newenvironment{paradoja}{%
  \begin{alertblock}{Paradoja de Exploración}%
}{%
  \end{alertblock}%
}

\newenvironment{motivacion}{%
  \begin{exampleblock}{Motivación: Ciencia de Datos en la Práctica}%
}{%
  \end{exampleblock}%
}

\newenvironment{retoclase}{%
  \begin{block}{Reto Rápido en Equipo}%
}{%
  \end{block}%
}


\title[Transformaciones No Lineales]{Sección 09.12: Transformaciones No Lineales y Regresión Polinomial}
\subtitle{El Cierre del Capítulo: Curvatura, Multicolinealidad Polinomial, y Síntesis}
\author[J. Castillo Colmenares]{Juliho Castillo Colmenares}
\institute{Tecnológico de Monterrey}
\date{\vspace{-1.2cm}}

\begin{document}

% Slide 1: Portada
\begin{frame}[plain]
  \titlepage
\end{frame}

% Slide 2: Hoja de Ruta
\begin{frame}{Hoja de Ruta --- Capítulo 09: Regresiones Lineales y Múltiples}
  \scriptsize
  ¿Dónde nos encontramos en el desarrollo modular del capítulo?
  \begin{itemize}\itemsep=0.02em
    \item \textbf{09.01} Correlación como Premisa de la Regresión
    \item \textbf{09.02} Introducción a la Regresión Lineal
    \item \textbf{09.03} Matemáticas de la Regresión: Derivación MCO y $R^2$
    \item \textbf{09.04} Regresión Lineal sobre Datos Simulados
    \item \textbf{09.05} Coeficientes Óptimos, Pruebas $t$/$F$ y RSE
    \item \textbf{09.06} Implementación en Python con \texttt{statsmodels}
    \item \textbf{09.07} Regresión Múltiple, Ecuación Normal y Ridge/Lasso
    \item \textbf{09.08} Validación de Modelos y $k$-fold Cross-Validation
    \item \textbf{09.09} Diagnóstico de Residuos y Multicolinealidad (VIF)
    \item \textbf{09.10} Regresión con \texttt{scikit-learn}
    \item \textbf{09.11} Variables Categóricas y Variables Muda
    \item \textbf{09.12} Transformaciones No Lineales y Regresión Polinomial \emph{(Hoy --- Cierre del Capítulo)}
  \end{itemize}
  \pause
  \vspace{-0.05cm}
  \begin{block}{Objetivo de la Sesión}
    Modelar relaciones curvas (no lineales en la variable original) reformulándolas como regresión lineal múltiple en variables expandidas; diagnosticar cuándo una transformación es necesaria; y cerrar el Capítulo 09 sintetizando su recorrido completo, desde la correlación hasta la regresión polinomial.
  \end{block}
\end{frame}

% Slide 3: Motivación
\begin{frame}{Cuando la Relación No Es una Línea Recta}
  \footnotesize
  \begin{motivacion}
    La potencia de un motor (\texttt{hp}) y su rendimiento (\texttt{mpg}) no siempre siguen una tendencia monótona simple: pueden formar una curva. Un modelo estrictamente lineal $\texttt{mpg}=c_0+c_1\,\texttt{hp}$ es incapaz de representar un cambio de dirección en la curva.
  \end{motivacion}

  \vspace{0.15cm}
  \pause
  \begin{itemize}\itemsep=0.1cm
    \item \textbf{Diagnóstico:} graficar la variable respuesta contra cada predictor revela la forma real de la relación. \pause
    \item \textbf{Solución:} agregar términos de orden superior ($X^2$, $X^3$, \ldots) o aplicar transformaciones (logaritmo, raíz cuadrada).
  \end{itemize}
\end{frame}

% Slide 4: Regresión polinomial como caso de regresión múltiple
\begin{frame}{La Regresión Polinomial Es Regresión Múltiple}
  \footnotesize
  \begin{block}{El Truco de la Variable Auxiliar}
    Definiendo $Z=X^2$ como una nueva columna calculada explícitamente, el modelo $Y=c_0+c_1X+c_2X^2+\varepsilon$ se reescribe exactamente como $Y=c_0+c_1X+c_2Z+\varepsilon$.
  \end{block}
  \pause

  \vspace{0.1cm}
  \begin{alertblock}{La Ecuación Normal No Cambia}
    Mínimos cuadrados solo exige linealidad \textbf{en los parámetros} $c_0,c_1,c_2$, no en $X$. La Ecuación Normal de la Sección 09.07, $\hat{\mathbf{c}}=(\mathbf{X}^T\mathbf{X})^{-1}\mathbf{X}^T\mathbf{Y}$, se aplica sin ninguna modificación.
  \end{alertblock}
\end{frame}

% Slide 5: Multicolinealidad polinomial
\begin{frame}{La Multicolinealidad Inherente de los Polinomios}
  \footnotesize
  \begin{block}{El Problema}
    Cuando $X$ está concentrada en un rango angosto, $X^2\approx X_0^2+2X_0(X-X_0)$ es casi una función lineal de $X$: $X$ y $X^2$ quedan altamente correlacionadas, inflando su VIF (Sección 09.09).
  \end{block}
  \pause

  \vspace{0.1cm}
  \begin{alertblock}{No Invalida las Predicciones}
    Las predicciones $\hat Y$ (que dependen de la combinación conjunta de $X$ y $X^2$) permanecen precisas; lo que se complica es la interpretación aislada de $c_1$ y $c_2$ por separado.
  \end{alertblock}
\end{frame}

% Slide 6: Lab Python (Bloque 1)
\begin{frame}[fragile]{Laboratorio en Python: Ajuste del Modelo Cuadrático}
  \scriptsize
  Expansión de la matriz de diseño y solución vía la Ecuación Normal (\texttt{09.12\_nonlinear\_polynomial\_regression.py}):
  {\lstset{basicstyle=\fontsize{3.2pt}{3.9pt}\selectfont\ttfamily,aboveskip=0.05em,belowskip=0.05em}
  \lstinputlisting[firstline=17, lastline=33]{../../code/09_regresiones/09.12_nonlinear_polynomial_regression.py}}
\end{frame}

% Slide 7: Lab Python (Bloque 2)
\begin{frame}[fragile]{Laboratorio en Python: Lineal vs. Cuadrático}
  \scriptsize
  Comparación de $R^2$ y Prueba $F$ Parcial sobre el término cuadrático:
  {\lstset{basicstyle=\fontsize{3.0pt}{3.7pt}\selectfont\ttfamily,aboveskip=0.05em,belowskip=0.05em}
  \lstinputlisting[firstline=38, lastline=59]{../../code/09_regresiones/09.12_nonlinear_polynomial_regression.py}}
\end{frame}

% Slide 8: Lab Python (Bloque 3)
\begin{frame}[fragile]{Laboratorio en Python: Multicolinealidad Polinomial}
  \scriptsize
  Correlación entre $X$ y $X^2$ según el rango de $X$:
  {\lstset{basicstyle=\fontsize{3.2pt}{3.9pt}\selectfont\ttfamily,aboveskip=0.05em,belowskip=0.05em}
  \lstinputlisting[firstline=62, lastline=75]{../../code/09_regresiones/09.12_nonlinear_polynomial_regression.py}}
\end{frame}

% Slide 9: Salida Terminal
\begin{frame}[fragile]{Laboratorio en Python: Salida en Terminal}
  \scriptsize
  Salida estándar generada al ejecutar el script del laboratorio en Python:
  \vspace{0.02cm}
  \begin{block}{Salida Estándar en Consola}
    \fontsize{3.2pt}{3.9pt}\selectfont
    \begin{verbatim}
=== Block 1: Quadratic Fit ===
Coefficients (c0,c1,c2): [64.2, -0.7634, 0.003657]; SSE=0.4571

=== Block 2: Linear vs. Quadratic ===
R^2 linear=0.0800, R^2 quadratic=0.9943, Partial F=320.0000

=== Block 3: Polynomial Multicollinearity ===
Narrow X in [95,105]: corr(X,X^2)=0.9999, VIF=5250.82
Wide X in [1,200]:    corr(X,X^2)=0.9673, VIF=15.56
    \end{verbatim}
  \end{block}
\end{frame}

% Slide 10A: Ejercicio Nivel 1 (Enunciado)

% Slide 10B: Ejercicio Nivel 1 (Resolución)

% Slide 11A: Ejercicio Nivel 2 (Enunciado)

% Slide 11B: Ejercicio Nivel 2 (Resolución)

% Slide 12A: Ejercicio Nivel 3 (Enunciado)
\begin{frame}{Ejercicio en Clase --- Nivel Analítico (1/2)}
  \footnotesize
  \begin{block}{Problema --- Polinomio como Regresión Lineal (Enunciado, Problema 9.20.7)}
    Demuestre que ajustar $Y=c_0+c_1X+c_2X^2+\varepsilon$ por mínimos cuadrados es formalmente idéntico a la regresión múltiple de la Sección 09.07 al definir $Z=X^2$, y explique por qué la Ecuación Normal se aplica sin modificación.
  \end{block}

  \vspace{0.1cm}
  \pause
  \begin{alertblock}{Estrategia}
    Distinga entre linealidad en los parámetros ($c_0,c_1,c_2$) y linealidad en el predictor original ($X$).
  \end{alertblock}
\end{frame}

% Slide 12B: Ejercicio Nivel 3 (Resolución)
\begin{frame}{Ejercicio en Clase --- Nivel Analítico (2/2)}
  \scriptsize
  Reescribimos el modelo con la variable auxiliar: \pause

  \vspace{0.05cm}
  \begin{block}{Cálculo}
    \scriptsize
    $$Y=c_0+c_1X+c_2X^2+\varepsilon \quad \xrightarrow{Z=X^2} \quad Y=c_0+c_1X+c_2Z+\varepsilon \quad \text{(lineal en } c_0,c_1,c_2\text{)}$$
  \end{block}

  \vspace{0.05cm}
  \pause
  \begin{alertblock}{Interpretación}
    Mínimos cuadrados solo requiere linealidad en los \textbf{parámetros}, nunca en las variables predictoras originales. Por eso la Ecuación Normal $(\mathbf{X}^T\mathbf{X})^{-1}\mathbf{X}^T\mathbf{Y}$ resuelve el problema sin ninguna generalización adicional.
  \end{alertblock}
\end{frame}

% Slide 13A: Ejercicio Nivel 4 (Enunciado)
\begin{frame}{Ejercicio en Clase --- Nivel Desafiante (1/2)}
  \footnotesize
  \begin{block}{Problema --- Multicolinealidad Polinomial (Enunciado, Problema 9.20.9)}
    Para $X\in[95,105]$ (rango angosto), demuestre que $X$ y $X^2$ están altamente correlacionadas, y explique --- conectando con la Sección 09.09 --- por qué esto infla el VIF sin invalidar las predicciones.
  \end{block}

  \vspace{0.1cm}
  \pause
  \begin{alertblock}{Estrategia}
    Escriba $X=X_0+\delta$ con $|\delta|$ pequeño y expanda $X^2$ en serie de Taylor.
  \end{alertblock}
\end{frame}

% Slide 13B: Ejercicio Nivel 4 (Resolución)
\begin{frame}{Ejercicio en Clase --- Nivel Desafiante (2/2)}
  \scriptsize
  Expandimos $X^2$ alrededor del centro del rango angosto: \pause

  \vspace{0.05cm}
  \begin{block}{Cálculo}
    \scriptsize
    $$X^2=(X_0+\delta)^2=X_0^2+2X_0\delta+\delta^2\approx X_0^2+2X_0\delta \quad (\delta^2 \text{ despreciable si } |\delta|\ll X_0)$$
  \end{block}

  \vspace{0.05cm}
  \pause
  \begin{alertblock}{Interpretación}
    $X^2$ se vuelve aproximadamente lineal en $X$ (pendiente $2X_0$), produciendo $r_{X,X^2}\approx1$ y VIF muy alto. Las predicciones $\hat Y$ (que dependen de $c_1X+c_2X^2$ conjuntamente) permanecen estables; solo se complica separar el "efecto de $X$" del "efecto de $X^2$" individualmente.
  \end{alertblock}
\end{frame}

% Slide 14: Cierre de Sección
\begin{frame}{Cierre de Sección: Transformaciones No Lineales}
  \begin{reflexion}
    La regresión polinomial demuestra la flexibilidad del marco de mínimos cuadrados: basta con expandir el espacio de predictores (agregar $X^2$, $X^3$, o transformaciones como $\log X$) para modelar curvatura, sin abandonar nunca la maquinaria lineal derivada en la Sección 09.07. La contrapartida es la multicolinealidad inherente entre potencias de la misma variable.
  \end{reflexion}

  \vspace{0.2cm}
  \pause
  \begin{block}{Lecciones Clave}
    \begin{itemize}\itemsep=0.08cm
      \item \textbf{Linealidad en los parámetros:} la verdadera condición de MCO, no la forma de la relación en $X$.
      \item \textbf{Prueba $F$ Parcial:} formaliza si un término de orden superior es realmente necesario.
      \item \textbf{Multicolinealidad polinomial:} VIF alto entre $X$ y $X^2$ no invalida las predicciones, solo la interpretación aislada.
    \end{itemize}
  \end{block}
\end{frame}

% Slide 15: Cierre del Capítulo
\begin{frame}{Cierre del Capítulo 09: Regresiones Lineales y Múltiples}
  \scriptsize
  \begin{retoclase}
    Con esta sección cerramos el Capítulo 09 completo: partimos de la \textbf{correlación} (09.01) como premisa conceptual, construimos la \textbf{regresión simple} desde su motivación hasta su implementación profesional (09.02-09.06), generalizamos a la \textbf{regresión múltiple} vía notación matricial y regularización (09.07), aprendimos a \textbf{validar} (09.08) y \textbf{auditar} (09.09) un modelo, lo implementamos con \textbf{herramientas de la industria} (09.10), y extendimos su alcance a predictores \textbf{categóricos} (09.11) y relaciones \textbf{no lineales} (09.12).
  \end{retoclase}

  \vspace{0.08cm}
  \pause
  \begin{alertblock}{El Hilo Conductor}
    Un único principio --- minimizar la suma de cuadrados de los residuos --- unifica cada una de estas doce secciones. El Capítulo 09 demuestra que la regresión lineal, lejos de ser una técnica limitada, es un marco extraordinariamente flexible sobre el cual se construyen prácticamente todos los métodos predictivos modernos de la ciencia de datos.
  \end{alertblock}
\end{frame}

\end{document}
